\section{Agent 能力的评估}

\subsection{为什么需要评估 Agent 能力}

\begin{importantbox}{评估的双重目的}
\begin{enumerate}
    \item \textbf{了解进展}：量化 Agent 在各类任务上的能力水平
    \item \textbf{预判风险}：识别接近危险阈值的能力维度，提前部署安全措施
\end{enumerate}
\end{importantbox}

\subsection{评估的维度}

\begin{itemize}
    \item \textbf{自主性}（Autonomy）：Agent 能否在无人监督下完成长周期任务？
    \item \textbf{工具使用}（Tool Use）：Agent 能多有效地使用外部工具（代码执行、搜索、API）？
    \item \textbf{说服力}（Persuasion）：Agent 能否影响人类决策？这涉及潜在的操纵风险。
    \item \textbf{CBRN 知识}：Agent 是否具备化学、生物、放射性和核领域的专业知识？
    \item \textbf{网络安全}：Agent 能否发现和利用软件漏洞？
\end{itemize}

\subsection{评估方法论}

\begin{itemize}
    \item 设计真实世界任务（非合成基准），观察 Agent 的端到端表现
    \item 人类基线对照：让领域专家完成同样任务，比较 Agent 表现
    \item 关注``尾部能力''：关注模型在极端情况下可能展现的能力
\end{itemize}

\begin{warningbox}{评估的根本困难}
评估是对称的——如果我们能评估一个能力，说明我们已经理解了这个能力。但最危险的可能恰恰是我们尚未识别的能力（unknown unknowns）。
\end{warningbox}

\subsection{本章小结}

Agent 能力评估需要覆盖自主性、工具使用、社会影响和专业知识等多个维度，且需要超越标准基准测试，关注真实世界表现和尾部风险。

